\documentclass[11pt,a4paper]{article}
\usepackage[margin=27mm]{geometry}\usepackage{fontspec,amsmath,amssymb,amsthm,unicode-math}
\setmainfont{STIXGeneral}\setmathfont{STIX Two Math}
\newtheorem{proposition}{Proposition}\newcommand{\gen}[1]{\langle#1\rangle}
\title{The converse for an endpoint central extension}
\author{Human Compression project}\date{Bounded correction to the specified author edition}
\begin{document}\maketitle
\begin{abstract}
A rank-three example distinguishes preservation of the intersection property under an endpoint central extension from its converse. The latter fails under the general hypothesis printed in Proposition 2.4 of the specified Cameron--Fernandes--Leemans author edition. This observation does not refute their stability theorem or the forward extension theorem used in our gluing applications.
\end{abstract}

\section{The exact source hypothesis}
The source is P. J. Cameron, M. E. Fernandes and D. Leemans, \emph{The number of string C-groups of high rank}, arXiv:2212.12723v2, dated 11 January 2023. Its definition on pp. 4--5 permits a string group generated by involutions when its ordered generating set consists of involutions with the distant commutations. An independent generating set and noncommuting adjacent generators are not imposed there.

In the notation of pp. 5--6, start with $G=\gen{\rho_0,\ldots,\rho_{r-1}}$ and a central involution $\tau$ in a supergroup, with $\tau\notin G$. Replacing $\rho_0$ by $\rho_0\tau$ gives the endpoint sesqui-extension $\Gamma^*$. Proposition 2.4 on p. 6 asserts, under just this hypothesis, that
\begin{quote}
``$\Gamma^*$ is a string C-group if and only if $\Gamma$ is a string C-group.''
\end{quote}
The assertion that the original is a C-group whenever the extension is a C-group is the direction considered here. We make no assertion about other editions or a later correction.

\section{An elementary abelian example}
\begin{proposition}
There is an endpoint sesqui-extension of a string group generated by three distinct involutions that is a C-group even though the original marking is not.
\end{proposition}
\begin{proof}
Let $G=\gen{a,b}\cong C_2^2$, with $a,b$ independent, and let $\tau$ generate a further direct factor of order two. Use the ordered marking
\[
 (\rho_0,\rho_1,\rho_2)=(a,b,ab).
\]
All three generators are distinct nonidentity involutions. They generate $G$ and commute, so in particular satisfy the string commutations. But
\[
 \gen{\rho_0}\cap\gen{\rho_1,\rho_2}
 =\gen{a}\ne 1,
\]
although the two index sets are disjoint. Thus the original marking fails the intersection property.

The endpoint extension has marking $(a\tau,b,ab)$. Its generators are independent: if
\[
 (a\tau)^{\epsilon_0}b^{\epsilon_1}(ab)^{\epsilon_2}=1,
 \qquad \epsilon_i\in\{0,1\},
\]
then the $\tau$-coordinate first gives $\epsilon_0=0$, the $a$-coordinate gives $\epsilon_2=0$, and the $b$-coordinate gives $\epsilon_1=0$. They therefore form a basis of $C_2^3$. A subgroup generated by any selected subset consists exactly of the vectors whose coordinates outside that subset vanish. Intersecting two such subgroups imposes both sets of zero coordinates, hence gives the subgroup on the common indices. This proves all intersection identities, not just the maximal-parabolic one.
\end{proof}

A faithful permutation realization is
\[
 a=(1\ 2)(3\ 4),\qquad b=(1\ 3)(2\ 4),\qquad
 ab=(1\ 4)(2\ 3),\qquad \tau=(5\ 6).
\]
The involution $\tau$ is outside $G$ and commutes with it, exactly as required in the source. The appended program checks these hypotheses and all sixty-four subgroup intersections for each of the two ordered triples. The written proof is independent of that finite computation.

\section{The invalid equality and the limited conclusion}
The last equality in the source's displayed proof chain identifies $G_0\cap G^*_{r-1}$ with $G_0\cap G_{r-1}$. In this example,
\[
 G_0=\gen{b,ab}=G,\qquad G_{2}=\gen{a,b}=G,
\]
whereas $G^*_{2}=\gen{a\tau,b}$ and
\[
 G_0\cap G^*_{2}=\gen{b},\qquad G_0\cap G_{2}=G.
\]
The central projection can therefore create an additional intersection. Equality of the two intersections is not automatic.

The forward statement, that twisting an endpoint of a C-marking by an independent central involution preserves the intersection property, is not contradicted. It is the direction retained in the admitted L010 application, where it is attributed to the earlier endpoint extension theorem. The present example does not satisfy the hypotheses of the source's rank-extension Theorem 6.4 as an input C-marking, and it does not disprove that theorem or the main stability result. Any use of the converse in a restricted permutation setting requires a further argument checking that setting. No new open-conjecture answer or complete L019 recovery is claimed here.
\end{document}
